\documentclass[11pt]{article}

\usepackage[T1]{fontenc}
\usepackage[utf8]{inputenc}
\usepackage{lmodern}
\usepackage[margin=1.15in, top=1in, bottom=1in]{geometry}
\usepackage{microtype}
\usepackage{booktabs}
\usepackage{enumitem}
\usepackage{titlesec}
\usepackage{fancyhdr}
\usepackage{xcolor}
\usepackage{parskip}
\usepackage{array}
\usepackage{setspace}

\definecolor{rule}{gray}{0.4}
\definecolor{muted}{gray}{0.35}

% Section formatting
\titleformat{\section}{\large\scshape\color{muted}}{}{0em}{}[\vspace{-4pt}\rule{\linewidth}{0.4pt}]
\titlespacing{\section}{0pt}{18pt}{8pt}

% No paragraph indent, small skip
\setlength{\parindent}{0pt}
\setlength{\parskip}{5pt}

% Compact lists
\setlist[itemize]{noitemsep, topsep=3pt, leftmargin=1.4em, label=\textendash}

% Header / footer
\pagestyle{fancy}
\fancyhf{}
\renewcommand{\headrulewidth}{0pt}
\fancyhead[L]{\small\color{muted}\textit{Applied Machine Learning for Business}}
\fancyhead[R]{\small\color{muted}\textit{Plaksha University, Monsoon 2026}}
\fancyfoot[C]{\small\color{muted}\thepage}

% Module heading command
\newcommand{\module}[2]{%
  \vspace{10pt}%
  {\large\textbf{#1\quad #2}}%
  \vspace{4pt}%
  \newline%
}

\begin{document}

% ── Title block ───────────────────────────────────────────────────────────────

\begin{center}
  {\LARGE\textbf{Applied Machine Learning for Business}}\\[6pt]
  {\large Plaksha University \quad|\quad Monsoon 2026}\\[10pt]
  \begin{tabular}{rl}
    \textit{Instructor} & Nikhil George \\
    \textit{Teaching Fellow} & Gajendra Dwivedi \\
    \textit{Office Hours} & By appointment\enspace---\enspace11\,am--5\,pm, Mon--Fri, Havells\,402 \\
  \end{tabular}
\end{center}

\vspace{6pt}
\noindent\rule{\linewidth}{0.6pt}
\vspace{4pt}

% ── Course Description ────────────────────────────────────────────────────────

\section*{Course Description}

Students arrive with machine learning foundations and programming fluency. This course builds on that preparation across a range of business situations where their technical competency in data science is called upon~--- from exploratory analysis and model building, to evaluating systems, to the economics and institutional consequences of how these systems behave. At each step the questions are precise: what is the objective function optimising, where does the metric break down, what does scale change. Some topics are named for a business domain~--- customer analytics, search and advertising~--- others for an algorithmic framework~--- networks, recommender systems~--- but in each, the work happens at the same intersection: the data, the algorithm, and the context in which it operates. Assignments span conceptual questions, data analysis, and programming.

\medskip
\textit{Prerequisites:} Prior coursework in machine learning required. Concurrent deep learning enrollment acceptable.

% ── Modules ───────────────────────────────────────────────────────────────────

\section*{Modules}

% --- Module 1 ---
\module{1.}{Customer Analytics}

A brief introduction to dimensionality reduction (PCA), customer segmentation, churn prediction, and customer lifetime value~--- including the limitations of each. A recurring theme: the choices made before the algorithm runs shape what the numbers report.

\smallskip
\textit{Strongly Encouraged}
\begin{itemize}
  \item McCarthy, ``A Playful Introduction to Customer Lifetime Value'' (2026)
  \item Ascarza, ``Retention Futility: Targeting High-Risk Customers Might Be Ineffective'' (2018) --- \textit{Journal of Marketing Research}
\end{itemize}

\textit{Technical Reference (optional)}
\begin{itemize}
  \item Abdi \& Williams, ``Principal Component Analysis'' (2010) --- \texttt{personal.utdallas.edu/\textasciitilde herve/abdi-awPCA2010.pdf}
\end{itemize}

\bigskip

% --- Module 2 ---
\module{2.}{Networks}

A node's importance is a function of its neighbors' importance~--- a recursive structure that degree centrality cannot capture. Eigenvector centrality, PageRank, community detection, and influence maximization: the algorithmic foundations of ranking and reach in a network.

\smallskip
\textit{Readings}
\begin{itemize}
  \item Barabasi, \textit{Network Science} (2016), Chapters 1--3 --- free at \texttt{networksciencebook.com}
  \item Page, Brin et al., ``The PageRank Citation Ranking'' (1999) --- Stanford Tech Report
  \item Haldane, ``Rethinking the Financial Network'' (2009) --- Bank of England speech
\end{itemize}

\bigskip

% --- Module 3 ---
\module{3.}{Evaluating Deployed Systems}

A screening algorithm can be evaluated only on cases it approved~--- a property called selective labels. This module introduces the problem, develops a technique for estimating true model performance under reasonable structural assumptions, and examines what standard metrics like AUC are actually measuring when the evaluation sample is not random. Bootstrap confidence intervals are introduced as a tool for quantifying uncertainty in performance estimates.

\smallskip
\textit{Strongly Encouraged}
\begin{itemize}
  \item Fawcett, ``An Introduction to ROC Analysis'' (2006) --- \textit{Pattern Recognition Letters}
\end{itemize}

\textit{Background Reference}
\begin{itemize}
  \item Kleinberg, Lakkaraju et al., ``Human Decisions and Machine Predictions'' (2018) --- \textit{QJE}; NBER working paper free
  \item Efron \& Hastie, \textit{Computer Age Statistical Inference}, Chapter~10 --- free PDF at Stanford
\end{itemize}

\bigskip

% --- Module 4 ---
\module{4.}{Search and Advertising}

Introduction to classical search and advertising. The module covers the auction mechanism that determines which ads are shown and at what price~--- including the generalized second-price auction. We then examine the cookie ecosystem and how it is changing. The attribution problem in advertising measurement comes next: an ad click that leads to a purchase is not straightforward to credit, and experiments are the primary tool for resolving it. The module closes with how deep learning has changed search and the frontier questions that arise when language models answer rather than rank.

\smallskip
\textit{Readings}
\begin{itemize}
  \item Varian, ``Position Auctions'' (2007) --- \textit{International Economic Review}; free on Varian's Berkeley page
  \item Edelman, Ostrovsky \& Schwarz, ``Internet Advertising and the Generalized Second-Price Auction'' (2007) --- \textit{AER}
  \item EFF, ``Behind the One-Way Mirror'' (2019) --- \texttt{eff.org}
  \item Kohavi et al., ``Controlled Experiments on the Web'' (2009) --- \textit{Data Mining and Knowledge Discovery}
\end{itemize}

\bigskip

% --- Module 5 ---
\module{5.}{Recommender Systems}

Users and items can be represented as vectors in a shared latent space~--- the dot product between them predicts preference. This is the underlying idea of recommender systems, from collaborative filtering to modern deep learning approaches. The module introduces the core concept, touches upon progress in the field, and then covers practical challenges: cold start, sparsity, and the fundamental difficulty of evaluating a recommender system without deployment.

\smallskip
\textit{Strongly Encouraged}
\begin{itemize}
  \item Netflix Technology Blog --- \texttt{netflixtechblog.medium.com}
\end{itemize}

\textit{Background Reference}
\begin{itemize}
  \item Koren, Bell \& Volinsky, ``Matrix Factorization Techniques for Recommender Systems'' (2009) --- \textit{IEEE Computer}
  \item Chaney, Stewart \& Engelhardt, ``How Algorithmic Confounding Increases Homogeneity'' (2018) --- arXiv:1710.11214
\end{itemize}

\bigskip

% --- Module 6 ---
\module{6.}{Economics of Modern AI}

An introduction to the economic and strategic questions raised by rapid progress in AI. The module opens with how machines learn representations~--- from word embeddings to recurrent networks to the transformer architecture~--- tracing the progression from count-based methods to the parallel, attention-based architecture that made large-scale training possible. From there, performance scales predictably with compute, data, and model size~--- an empirical regularity with known corrections and contested limits. The module covers the historical argument that scale beats specialization (Sutton), what the scaling laws actually claim and where they break (Kaplan, Hoffmann), and what it means that frontier AI is beginning to produce new results in mathematics and science (Bubeck et al.).

\smallskip
\textit{Strongly Encouraged}
\begin{itemize}
  \item Sutton, ``The Bitter Lesson'' (2019) --- \texttt{incompleteideas.net}
  \item Kaplan et al., ``Scaling Laws for Neural Language Models'' (2020) --- arXiv:2001.08361
  \item Bubeck et al., ``Early Science Acceleration with GPT-5'' (2025)
  \item Eric Jang, ``Building AlphaGo from Scratch with Modern AI Tools'' [video] --- \texttt{youtube.com/watch?v=X\_ZVSPcZhtw}
\end{itemize}

\textit{Background Reference}
\begin{itemize}
  \item Wigner, ``The Unreasonable Effectiveness of Mathematics in the Natural Sciences'' (1960)
  \item Hoffmann et al., ``Training Compute-Optimal Large Language Models'' (2022) --- arXiv:2203.15556
  \item Amodei, ``Machines of Loving Grace'' (2024) --- \texttt{darioamodei.com}
  \item Thompson et al., ``The Computational Limits of Deep Learning'' (2020) --- arXiv:2007.05558
\end{itemize}

\bigskip

% --- Module 7 (Time Permitting) ---
\module{7.}{Time Permitting}

How AI progress is measured~--- and why that is hard: Goodhart's law applied to benchmarks, contamination, and the claim that emergent capabilities are an artifact of discontinuous metrics rather than discontinuous capability. The module closes with the immediate business consequence: a two-hundred-billion-dollar advertising model built on search links is under structural threat from systems that answer rather than rank.

% ── Assessment ────────────────────────────────────────────────────────────────

\section*{Assessment}

\begin{center}
\renewcommand{\arraystretch}{1.35}
\begin{tabular}{lc}
  \toprule
  \textbf{Component} & \textbf{Weight} \\
  \midrule
  Team Assignments (3)       & 35\% \\
  In-Class Tests (2)         & 30\% \\
  Individual Assignments (2) & 25\% \\
  Participation              & 10\% \\
  \bottomrule
\end{tabular}
\end{center}

\medskip
Class tests are conducted in-class and drawn from assignment questions. Attendance at or above 80\% of sessions earns full participation credit; below that threshold earns zero.

\end{document}
